ФИО: Есаулова Наталья Андреевна
Тема ВКР: "Разработка мультиагентной конфигурации для оптимизации решения исследовательских и научных задач на базе OpenClaw"
Задание ВКР:
\begin{itemize}
  \item Проведение анализа подходов к применению мультиагентных систем для решения научных и исследовательских задач.
  \item Проведение анализа архитектуры и функциональных возможностей OpenClaw.
  \item Проектирование расширенной мультиагентной конфигурации OpenClaw, направленной на повышение эффективности решения научных задач.
  \item Реализация модулей и механизмов их взаимодействия в данной конфигурации.
  \item Разработка MVP-сервиса для тестирования и демонстрации работы системы.
  \item Проведение тестирования системы и оценка ее эффективности.
\end{itemize}

Аннотация ВКР:

Цель исследования

Целью работы является разработка и проверка архитектурной доработки агентной платформы на базе OpenClaw агентов ClawdLab для повышения эффективности выполнения научных задач. 

Задачи, решаемые в ВКР

В работе решаются следующие задачи: анализ предметной области и современного состояния агентных систем для научных задач; анализ архитектуры OpenClaw и ClawdLab, их возможностей и ограничений; формулирование архитектурной гипотезы о применении структурированной памяти и управляемой цикличности выполнения задач; проектирование экспериментальной схемы для сравнения базовой и улучшенной конфигураций; разработка MVP-контура ClawdLab на собственной инфраструктуре; реализация доработкок; оценка результатов в различных конфигурациях системы; подведение итогов.

Краткая характеристика полученных результатов

В рамках работы был проведен обзор существующих подходов к применению мультиагентных LLM-систем в научных задачах и выделены архитектурные ограничения ClawdLab, связанные с потерей контекста, недостаточной трассируемостью артефактов и неформальным характером цикличности выполнения задач. Был спроектирован экспериментальный контур для проверки доработок на задачах анализа и глубокого исследования в рамках определенного исследовательского вопроса. Также был развернут MVP-сервис ClawdLab и агенты OpenClaw и выполнены необходимые доработки системы для ее тестирования.

Основным практическим результатом стала реализация двух доработок: структурированной памяти и управляемой цикличности выполнения задач. Для памяти была предложена и реализована трехуровневая модель, включающая рабочий, эпизодический и канонический уровни, а также сущности для хранения записей памяти, ссылок на источники и продвижения объектов памяти между уровнями. Для цикличности была спроектирована и добавлена логика фиксированных правил повторного выполнения задач, позволяющая связывать новую попытку с критикой, причиной повторного запуска и условиями остановки. Вместе эти доработки позволяют сохранять происхождение выводов, связывать результаты задач с источниками и артефактами, а также делать процесс выполнения научных задач более прозрачным и воспроизводимым.
